\documentclass[11pt,a4paper]{article}
\usepackage[utf8]{inputenc}
\usepackage[italian]{babel}
\usepackage{geometry}
\geometry{top=2.5cm,bottom=2.5cm,left=2.5cm,right=2.5cm}
\usepackage{amsmath,amssymb}
\usepackage{tcolorbox}
\tcbuselibrary{skins,breakable}
\usepackage{booktabs}
\usepackage{tabularx}
\usepackage{enumitem}
\usepackage{xcolor}
\usepackage{listings}
\usepackage{hyperref}
\hypersetup{
    colorlinks=true,
    linkcolor=blue!70!black,
    citecolor=blue!70!black,
    urlcolor=teal!70!black
}

\lstset{
    basicstyle=\ttfamily\small,
    breaklines=true,
    frame=single,
    backgroundcolor=\color{gray!10},
    keywordstyle=\color{blue}\bfseries,
    commentstyle=\color{green!50!black}\itshape,
    stringstyle=\color{red!70!black}
}

\newtcolorbox{questionbox}[1][]{
    colback=gray!5!white,
    colframe=gray!60!black,
    fonttitle=\bfseries,
    title={#1},
    breakable
}

\newtcolorbox{solutionbox}[1][]{
    colback=blue!5!white,
    colframe=blue!70!black,
    fonttitle=\bfseries,
    title={#1},
    breakable
}

\newtcolorbox{altbox}[1][]{
    colback=orange!5!white,
    colframe=orange!80!black,
    fonttitle=\bfseries,
    title={#1},
    breakable
}

\begin{document}

\title{\textbf{Sicurezza dei Sistemi Informativi} \\ \large Appello d'Esame del 25 Giugno 2024 \\ \small Anno Accademico 2023-2024}
\author{Politecnico di Torino --- Docente: Prof. Cataldo Basile}
\date{2024-06-25}
\maketitle


\section{Testo Ufficiale dell'Esame}

\begin{questionbox}[Domanda 1: Domande a Scelta Multipla (Vulnerabilit\`a e Protezione)]
Valutare la correttezza delle seguenti affermazioni:
\begin{enumerate}[label=\alph*.]
    \item Una vulnerabilit\`a di tipo buffer overflow sullo stack permette a un attaccante di sovrascrivere il frame pointer (Saved EBP/RBP) e l'indirizzo di ritorno (Saved EIP/RIP).
    \item La tecnica ASLR (Address Space Layout Randomization) rende completamente impossibili gli attacchi basati su Return-Oriented Programming (ROP).
    \item Gli stack canaries sono valori casuali inseriti prima dell'indirizzo di ritorno che vengono verificati prima dell'uscita della funzione per rilevare la corruzione dello stack.
    \item Il meccanismo Data Execution Prevention (DEP / NX bit) impedisce l'esecuzione di codice presente nelle sezioni stack e heap della memoria.
    \item Il processo di Reverse Engineering su codice compilato produce sempre come output il codice sorgente originale identico a quello scritto dallo sviluppatore.
\end{enumerate}
\end{questionbox}

\begin{questionbox}[Domanda 2: Phishing Avanzato, Pretexting e Debolezze Umane]
Spiegare cosa sono il \textbf{Phishing} e il \textbf{Pretexting}, evidenziando le differenze tra le due tecniche.
Descrivere le varianti principali del phishing:
\begin{itemize}
    \item Spear Phishing;
    \item Whaling;
    \item Vishing e Smishing.
\end{itemize}
Indicare quali debolezze psicologiche e cognitive umane vengono sfruttate (es. autorit\`a, urgenza, timore di sanzioni, avidit\`a) e proporre un piano di contromisure integrato (tecnologiche e formative).
\end{questionbox}

\begin{questionbox}[Domanda 3: Esercizio di Progettazione Architetturale Screened Subnet]
Un'organizzazione deve dividere la propria infrastruttura in due livelli di sicurezza distinti:
\begin{itemize}
    \item Rete a livello di sicurezza pi\`u basso (pu\`o accedere direttamente ad Internet);
    \item Rete a livello di sicurezza pi\`u elevato (separata da Internet e accessibile solo attraverso la rete a livello pi\`u basso tramite rigide regole di autorizzazione);
    \item Due server aziendali che erogano servizi critici devono essere raggiungibili anche da Internet.
\end{itemize}
Progettare la topologia di rete, specificare il posizionamento dei dispositivi di sicurezza (firewall, router, DMZ) e compilare la matrice delle regole di autorizzazione.
\end{questionbox}

\newpage
\section{Soluzione Dettagliata e Guida Risolutiva}

\begin{solutionbox}[Soluzione Domanda 1: Valutazione Crocette]
\begin{itemize}
    \item \textbf{a. VERA}: Un buffer overflow classico su array allocati sullo stack consente, scrivendo oltre la dimensione consentita, di sovrascrivere variabili locali, il Saved Frame Pointer e il Return Address, deviando il flusso di esecuzione.
    \item \textbf{b. FALSA}: L'ASLR randomizza gli indirizzi base delle librerie e dello stack, ma NON rende impossibili gli attacchi ROP; un attaccante pu\`o aggirarla tramite \textit{Information Disclosure / Memory Leak} per calcolare l'indirizzo base a runtime, oppure effettuare attacchi di forza bruta (in sistemi a 32 bit).
    \item \textbf{c. VERA}: Lo stack canary \`e una sentinella posta tra le variabili locali e il return address; se l'epilogo della funzione rileva un valore alterato, il programma invoca immediatamente \texttt{\_\_stack\_chk\_fail} terminando l'esecuzione.
    \item \textbf{d. VERA}: Il flag NX/XD (\textit{No-Execute}) a livello di paginazione di memoria hardware impedisce alla CPU di eseguire istruzioni fetchate da pagine marcate come dati (stack e heap), forzando gli attaccanti a passare ad attacchi di tipo code-reuse (Return-to-libc, ROP).
    \item \textbf{e. FALSA}: Il compilatore elimina nomi di variabili locali, commenti, tipi complessi e macro; i decompilatori moderni (Ghidra, Hex-Rays) ricostruiscono una rappresentazione ad alto livello approssimata, non il sorgente identico originario.
\end{itemize}
\end{solutionbox}

\begin{solutionbox}[Soluzione Domanda 2: Phishing, Pretexting e Psicologia dell'Ingegneria Sociale]
\textbf{Phishing vs Pretexting:}
\begin{itemize}
    \item \textbf{Phishing}: Attacco condotto tipicamente in massa via email o messaggistica che induce le vittime a compiere un'azione immediata (cliccare su un link clone o aprire un malware allegato).
    \item \textbf{Pretexting}: Creazione di uno scenario fittizio credibile e articolato (un "pretesto") in cui l'attaccante impersona una figura autorevole o di fiducia (es. auditor esterno, supporto IT, personale delle forze dell'ordine) per instaurare una relazione e farsi consegnare informazioni riservate nel corso di pi\`u interazioni.
\end{itemize}

\textbf{Varianti del Phishing:}
\begin{itemize}
    \item \textit{Spear Phishing}: Attacco altamente mirato contro un individuo o reparto specifico, personalizzato raccogliendo informazioni preliminari OSINT.
    \item \textit{Whaling}: Spear phishing diretto contro figure apicali aziendali (CEO, CFO, membri del CDA) con l'obiettivo di frodi finanziarie ingenti (es. CEO Fraud / Business Email Compromise).
    \item \textit{Vishing (Voice Phishing)}: Condotto telefonicamente, spesso oggi integrato con audio deepfake vocale generativo.
    \item \textit{Smishing}: Condotto tramite SMS contenenti link ingannevoli per intercettare credenziali o codici OTP SMS.
\end{itemize}

\textbf{Debolezze Psicologiche Sfruttate:} Senso di urgenza artificioso, timore di sanzioni o blocco conto, principio di autorit\`a, compiacenza e curiosit\`.

\textbf{Contromisure:}
\begin{itemize}
    \item \textit{Tecniche}: Adozione di chiavi hardware FIDO2 (resistenti al phishing), autenticazione email (DMARC su policy reject), sandboxing avanzato delle email in ingresso.
    \item \textit{Formative}: Campagne periodiche di simulazione phishing con debriefing immediato e cultura aziendale priva di colpevolizzazione (\textit{no-blame culture}).
\end{itemize}
\end{solutionbox}

\begin{solutionbox}[Soluzione Domanda 3: Topologia a Due Livelli di Sicurezza e DMZ]
\textbf{Architettura Progettata (Screened Subnet a Due Livelli di Sicurezza):}
Si adotta una struttura perimetrale a zone gerarchiche con due linee di difesa firewall distinte (Two-Tier Architecture):
\begin{center}
\begin{verbatim}
                         +--------------------+
                         |      INTERNET      | (Livello 0: Untrusted)
                         +---------+----------+
                                   | (Traffico Pubblico TCP 80/443)
                       +-----------v-----------+
                       |  FIREWALL ESTERNO L4  |
                       +-----------+-----------+
                                   |
         +-------------------------v-------------------------+
         |      LIVELLO 1: DMZ SCREENED SUBNET (192.168.1.0/24)|
         |                                                   |
         |   [Server Aziendale 1 - Web]  192.168.1.10        |
         |   [Server Aziendale 2 - API]  192.168.1.20        |
         |   [WAF / Reverse Proxy]       192.168.1.5         |
         +-------------------------+-------------------------+
                                   |
                       +-----------v-----------+
                       |  FIREWALL INTERNO L4  |
                       +-----------+-----------+
                                   |
         +-------------------------v-------------------------+
         |     LIVELLO 2: HIGH SECURITY LAN (10.10.0.0/16)   |
         |                                                   |
         |   [Database Master & Storage] 10.10.1.50          |
         |   [Workstation Riservate]     10.10.2.0/24        |
         |   [Proxy Uscita Aziendale]    10.10.1.10          |
         +---------------------------------------------------+
\end{verbatim}
\end{center}

\textbf{Subnetting, Instradamento e Flussi di Traffico:}
\begin{itemize}
    \item \textbf{Livello 0 (Untrusted / Internet)}: Accesso pubblico consentito esclusivamente verso il WAF/Server 1 e Server 2 in DMZ su porte TCP 80 e 443 tramite Destination NAT (DNAT). Nessun pacchetto esterno pu\`o mai raggiungere il Livello 2.
    \item \textbf{Livello 1 (DMZ, \texttt{192.168.1.0/24})}: Gateway \texttt{192.168.1.1} su Firewall Esterno e \texttt{192.168.1.2} su Firewall Interno. Ospita i server critici esposti. In caso di compromissione di un server in DMZ, il Firewall Interno blocca rigorosamente qualsiasi movimento laterale verso la LAN interna.
    \item \textbf{Livello 2 (High Security LAN, \texttt{10.10.0.0/16})}: Gateway \texttt{10.10.0.1}. Rete completamente schermata e isolata. Pu\`o raggiungere i server in DMZ solo per manutenzione su porte specifiche (SSH TCP 22 / SQL TCP 5432). L'uscita verso Internet \`e mediata e filtrata da un proxy aziendale autenticato (\texttt{10.10.1.10}).
\end{itemize}

\textbf{Matrice delle Regole di Autorizzazione:}
\begin{center}
\small
\begin{tabularx}{\textwidth}{llllX}
\toprule
\textbf{Origine} & \textbf{Destinazione} & \textbf{Protocollo/Porta} & \textbf{Azione} & \textbf{Motivazione} \\
\midrule
Internet & DMZ (Server 1, 2) & TCP 80, 443 & \textbf{PERMIT} & Esposizione dei due servizi pubblici aziendali \\
Internet & High Security LAN & Qualsiasi & \textbf{DENY} & Isolamento assoluto della rete interna da Internet \\
High Security LAN & DMZ & Porte dedicate & \textbf{PERMIT} & Amministrazione e aggiornamento dei server in DMZ \\
High Security LAN & Internet & TCP 80, 443 & \textbf{PERMIT} & Uscita web mediata da proxy aziendale \\
DMZ & High Security LAN & Qualsiasi & \textbf{DENY} & Previene movimenti laterali in caso di compromissione server DMZ \\
\bottomrule
\end{tabularx}
\end{center}
\end{solutionbox}

\begin{altbox}[Alternative Implementation Note: Rust per Software Anti-Exploit]
In Rust, la gestione della memoria \`e verificata staticamente dal \textit{borrow checker} del compilatore:
\begin{itemize}
    \item Le vulnerabilit\`a di stack overflow, buffer overflow e use-after-free sono \textbf{impossibili a tempo di compilazione} nel codice Safe Rust.
    \item Per la programmazione concorrente sicura senza race condition, si utilizzano tipi RAII standard come \texttt{std::sync::Mutex} e \texttt{std::sync::mpsc}.
\end{itemize}
\end{altbox}


\end{document}
